\section{Fondements scripturaires et stoïciens de la pensée politique ambrosienne}

\subsection*{Introduction}

\subsection{L’exemple des rois bibliques pour sortir des conceptions théoriques}

\subsubsection{Le bon souverain défini par l'Ancien Testament}

\subsubsection{Légitimer son pouvoir par l'implication religieuse}

Au-delà de la simple théorisation de l'origine divine du pouvoir, l'utilisation des figures bibliques permet à Ambroise de Milan de repenser entièrement les fondements de la légitimité impériale, ainsi que la grandeur des actions menées par le souverain contemporain. Comme nous l'avons abordé dans la première division de ce chapitre, l'objectif de l'évêque est de laisser de côté les conceptions purement philosophiques ou théoriques pour ancrer sa réflexion dans une praxis politique. Il s'agit de comprendre ce que recherche véritablement Ambroise chez les empereurs chrétiens afin de faire coïncider la réalité du pouvoir impérial avec son image idéale de la monarchie chrétienne. Un axe fondamental se dégage alors de cette démarche : la légitimité de l'empereur n'est plus un acquis institutionnel figé, mais une donnée proportionnelle à son engagement religieux. Plus l'implication du souverain dans la défense et la promotion de la foi chrétienne — et plus précisément de l'orthodoxie nicéenne — est élevée, plus la légitimité de son pouvoir s'en trouve renforcée.

\bigskip

Cette conditionnalité de la légitimité impériale s'exprime avec une clarté remarquable dans la correspondance d'Ambroise, notamment lorsqu'il s'adresse au jeune Valentinien II dans le contexte brûlant de la querelle de l'Autel de la Victoire. L'évêque de Milan n'hésite pas à redéfinir la hiérarchie du monde romain pour soumettre l'empereur à un devoir de service absolu envers la divinité :

\bigskip

\begin{quote}
«~Tous ceux qui sont sous la domination de Rome sont enrôlés pour vous servir, vous, les empereurs et rois de la terre ; ainsi, vous êtes vous-mêmes enrôlés pour servir le Dieu Tout-Puissant et notre sainte foi. Car la sécurité ne peut être mise en péril que si chacun n'est pas un véritable adorateur du vrai Dieu, le Dieu des chrétiens, qui gouverne toutes choses. Lui seul est le vrai Dieu et doit être adoré par l'esprit le plus profond. Quant à tous les dieux des païens, ce ne sont que des idoles, comme l'Écriture le dit.\footnote{Ambroise de Milan, \textit{Epistolae}, 72, 1. [Ajouter ici la citation latine et la balise cite si besoin]}~»
\end{quote}

\bigskip

Ce passage est fondamental pour saisir le basculement idéologique opéré par Ambroise. En utilisant le vocabulaire de l'engagement militaire («~enrôlés pour servir~», qui renvoie au concept de \latin{militia}), il dresse une pyramide du pouvoir strictement descendante. Si les populations de l'Empire doivent obéissance à l'empereur, ce dernier est à son tour le "soldat" du Christ. L'autorité impériale ne se justifie plus par elle-même, ni même par le Sénat ou les légions, mais par le service actif rendu à la «~sainte foi~». Plus encore, Ambroise opère un chantage politique particulièrement habile en liant la sécurité même de l'État (\latin{salus rei publicae}) à la dévotion exclusive du souverain. La \latin{Pax Romana} face aux périls extérieurs et intérieurs ne dépend plus de la protection des anciens dieux de Rome, mais de la ferveur nicéenne de l'empereur. Si le souverain se montre tolérant envers le polythéisme, il met l'Empire en péril. L'intransigeance religieuse devient ainsi, dans le discours ambrosien, une question de salut public et de sécurité nationale.

\bigskip

Pour illustrer et justifier ce besoin d'un souverain agissant concrètement contre les cultes païens, Ambroise recourt à une figure biblique majeure de l'Ancien Testament : celle du roi Josias. Reconnu par la tradition scripturaire comme le roi réformateur par excellence, Josias est celui qui a combattu le culte des idoles, purifié le Temple et restauré l'Alliance pour célébrer le Dieu unique. Cette figure permet à Ambroise de dresser un parallèle direct avec l'action politique de l'empereur Théodose. Dans l'oraison funèbre prononcée à la mort de ce dernier, l'évêque célèbre la lutte du souverain contre le polythéisme en fusionnant la mémoire impériale et la mémoire biblique :

\bigskip

\begin{quote}
«~Celui qui, sur cette terre, aura bien célébré la Pâque du Seigneur vivra dans la lumière sans fin. Or, qui l’a célébrée avec plus d’éclat que celui qui a refoulé les erreurs sacrilèges, fermé les temples, détruit les statues ? C’est pour cette raison que le roi Josias a été placé avant les plus grands.\footnote{Ambroise de Milan, \textit{De Obitu Theodosii}, 38. [Ajouter balise cite]}~»
\end{quote}

\bigskip

Cette comparaison n'est en rien fortuite. En assimilant Théodose à Josias, Ambroise sacralise la répression du paganisme. Fermer les temples et détruire les statues ne sont plus perçus comme des actes de violence tyrannique ou de rupture avec la tradition romaine ancestrale (\latin{mos maiorum}), mais comme l'accomplissement d'un devoir royal sacré. La légitimité de Théodose, et sa supériorité sur les "plus grands" souverains du passé, repose exclusivement sur son action iconoclaste et sa purification de l'Empire. Ambroise insiste sur cette réussite quelques paragraphes plus haut dans la même oraison funèbre, martelant que c'est la foi de l'empereur qui a permis ce bouleversement sociétal :

\bigskip

\begin{quote}
«~[...] sa foi a en effet fait disparaître tous les cultes des idoles, aboli toutes leurs cérémonies, lui a regretté que le pardon qu’il avait accordé à ceux qui avaient péché contre lui ait été perdu et qu’ils aient refusé sa grâce.\footnote{Ambroise de Milan, \textit{De Obitu Theodosii}, 4. [Ajouter balise cite]}~»
\end{quote}

\bigskip

À deux reprises dans le \latin{De Obitu Theodosii}, la question de la lutte contre le polythéisme est ainsi érigée en vertu suprême du gouvernement chrétien. La mention de la destruction des idoles par un souverain en exercice résonne profondément avec le Livre des Rois, mais elle fait surtout écho aux actions politiques et législatives bien réelles menées par Théodose. En effet, l'idéal théorique du roi Josias prôné par Ambroise trouve une traduction juridique implacable dans la politique religieuse de l'Empire, marquée par une gradation sévère des interdictions entre 381 et 392, visant à asphyxier la religion romaine.

\bigskip

La législation impériale, conservée dans le \latin{Codex Theodosianus}, illustre parfaitement cette mise en conformité du pouvoir avec les attentes ecclésiastiques. Dans un premier temps, les mesures s'attaquent aux manifestations les plus publiques du paganisme, par le biais d'une interdiction renouvelée des sacrifices sanglants (\latin{CTh} 16, 10, 8). L'étau législatif se resserre ensuite de manière ciblée, condamnant sévèrement les rassemblements dans les temples et la vénération des statues, notamment dans les bastions symboliques du paganisme que sont la ville de Rome et la province d'Égypte (\latin{CTh} 16, 10, 10 et 16, 10, 11). Enfin, l'aboutissement de cette politique, qui fait définitivement de Théodose le "nouveau Josias" célébré par Ambroise, survient en 392, alors qu'il règne en tant que seul empereur légitime de l'Orient et de l'Occident. L'édit de Constantinople (\latin{CTh} 16, 10, 12) constitue la mesure la plus radicale de l'Antiquité tardive : il proscrit absolument tout rassemblement, interdit les sacrifices envers les idoles et les dieux de la religion romaine, y compris dans la sphère domestique et privée (culte des Lares et des Pénates), et va jusqu'à valoriser la dénonciation de ces pratiques, désormais assimilées à des actes de lèse-majesté.

\bigskip

Ainsi, la pensée politique d'Ambroise ne se contente pas de fournir une caution morale au pouvoir en place. En utilisant le précédent biblique de Josias et en affirmant que l'empereur est "enrôlé" pour le service du vrai Dieu, l'évêque lie indéfectiblement la grandeur temporelle de Rome à la destruction du paganisme. La légitimité impériale est donc redéfinie : l'empereur n'est légitime que s'il est le bras séculier de la foi chrétienne, transformant la contrainte religieuse imposée par la loi romaine en un acte de piété digne des plus grands rois d'Israël.
